{Note: The trigger and reconstruction efficiencies discussed here are
  identical to those used in HION-2020-10~\cite{HION-2020-10}.
In a quick review this section may be skipped.
The \pp\ reconstruction and trigger efficiencies are taken from
  the publication HION-2019-58~\cite{HION-2019-58} and are not repeated here.
}

In constructing the pair distributions, a correction is
  applied accounting for the inefficiency arising from the trigger
  selection and the inefficiency of offline muon reconstruction.
The correction is applied by applying a weight to each dimuon pair,
  $w$, defined as :
\begin{eqnarray}
w^{-1}\equiv \varepsilon_{\mathrm{trig}} (\vec{p}_{a},\vec{p}_{b})
\varepsilon_{\mathrm{reco}}(\vec{p}_{a})  \varepsilon_{\mathrm{reco}}(\vec{p}_{b})
\label{eq:net_eff_correction}
\end{eqnarray}
where $\varepsilon_{\mathrm{trig}}$ is the efficiency
  of the dimuon trigger,  $\varepsilon_{\mathrm{reco}}$ is the
  single muon reconstruction efficiency. 
These are disscussed in sections~\ref{sec:trig_eff}--\ref{sec:reco_eff} below.


\subsection{Triggers efficiencies\label{sec:trig_eff}}

As mentioned before in Section~\ref{sec:data}, events are selected using a combination
  of the following di-muon triggers:
\begin{itemize}
\item \verb=HLT_mu4_mu4noL1= : Main trigger that accounts for most of the statistics.
\item \verb=HLT_2mu4= : Adds further few percent statistics.
\end{itemize}


It is assumed that the trigger efficiency for \verb=HLT_2mu4= factorizes into the
  product of the single muon trigger efficiencies :
\begin{eqnarray}
\varepsilon_{\mathrm{trig}}^{\mathrm{2mu4}}
(\vec{p}_{a},\vec{p}_{b}) \equiv
\varepsilon_{\mathrm{trig}}^{\mathrm{mu4}}   (\vec{p}_{a})
\varepsilon_{\mathrm{trig}}^{\mathrm{mu4}}   (\vec{p}_{b}),
\label{eq:2mu3_eff}
\end{eqnarray}
where $\varepsilon_{\mathrm{trig}}^{\mathrm{mu4}} $ is the single muon
  trigger-efficiency from matching an offline-muon with 
  a HLT muon with $\Delta R < 0.1$.

While for the \verb=HLT_mu4_mu4noL1= trigger the efficiency is given by
\begin{align}
\varepsilon_{\mathrm{trig}}^{\mathrm{mu4\_mu4noL1}}
(\vec{p}_{a},\vec{p}_{b}) \equiv&
 \varepsilon_{\mathrm{trig}}^{\mathrm{mu4}} (\vec{p}_{a})\ \times\  \varepsilon_{\mathrm{trig}}^{\mathrm{mu4noL1}} (\vec{p}_{b})\ +\ 
 \varepsilon_{\mathrm{trig}}^{\mathrm{mu4}} (\vec{p}_{b})\ \times\  \varepsilon_{\mathrm{trig}}^{\mathrm{mu4noL1}} (\vec{p}_{a})\nonumber  \\
 -\ & 
(\varepsilon_{\mathrm{trig}}^{\mathrm{mu4}} (\vec{p}_{a}) \cap \varepsilon_{\mathrm{trig}}^{\mathrm{mu4noL1}} (\vec{p}_{a}))\ \times\ 
(\varepsilon_{\mathrm{trig}}^{\mathrm{mu4}} (\vec{p}_{b}) \cap \varepsilon_{\mathrm{trig}}^{\mathrm{mu4noL1}} (\vec{p}_{b})),
\label{eq:mu4_mu4_nol1_eff}
\end{align}
where, in the second line, $(\varepsilon_{\mathrm{trig}}^{\mathrm{mu4}} (\vec{p}_{a}) \cap \varepsilon_{\mathrm{trig}}^{\mathrm{mu4noL1}} (\vec{p}_{a}))$
  is the probability for a muon to simultaneously fire the  \verb=HLT_mu4= and \verb=HLT_mu4noL1= chains.
Eq.~\eqref{eq:mu4_mu4_nol1_eff} states that the net efficiency for \verb=HLT_mu4_mu4noL1=
  is the union of the efficiency for one muon to fire the \verb=HLT_mu4= chain,  
  and the other muon to fire the \verb=HLT_mu4noL1= chain, while accounting for the double counting
  in the case both muons fire both chains.


The \verb=HLT_mu4_mu4noL1= trigger is the primary trigger used in this analysis
  and accounts for most of the statistics.
In principle, it would be expected that events selected by the \verb=HLT_2mu4= 
  trigger would be a subset of the events selected by the \verb=HLT_mu4_mu4noL1= trigger.
However, the HLT hypos for the \verb=HLT_mu4= and \verb=HLT_mu4noL1= 
  chains are different, resulting in \verb=HLT_mu4_mu4noL1= not being 
  a superset of the \verb=HLT_2mu4= trigger.
Additionally, because the \verb=HLT_mu4_mu4noL1= was prescaled in a few runs,
  while the \verb=HLT_2mu4= trigger was not, including the \verb=HLT_2mu4=
  trigger gains some additional statistics.
The following strategy is used for combining the \verb=HLT_mu4_mu4noL1= 
  and \verb=HLT_2mu4= triggers:
\begin{enumerate}
  \item Whenever a pair matches the \verb=HLT_mu4_mu4noL1=, the efficiency correction for \verb=HLT_mu4_mu4noL1= (Eq.~\eqref{eq:mu4_mu4_nol1_eff}) is used (and step 2 below is skipped).
  \item For pairs that do not match the  \verb=HLT_mu4_mu4noL1=, but match the \verb=HLT_2mu4= trigger, it is checked if \verb=HLT_mu4_mu4noL1= was prescaled out in that event (this information is stored in the xAOD):
  \begin{enumerate}[label=\alph*.]
    \item If \verb=HLT_mu4_mu4noL1= was prescaled out in this event, then this pair was selected by the \verb=HLT_2mu4 trigger=, and the efficiency correction for \verb=HLT_2mu4= (Eq.~\eqref{eq:2mu3_eff}) is applied.
    \item If \verb=HLT_mu4_mu4noL1= was not prescaled for this event, then this pair is not used (This is considered as a pair lost because of the inefficiency of the \verb=HLT_mu4_mu4noL1= trigger.
  \end{enumerate}
\end{enumerate}
Put another way, when the \verb=HLT_mu4_mu4noL1= trigger is not prescaled,
  its decision is used (without looking at the \verb=HLT_2mu4= decision.).
When the \verb=HLT_mu4_mu4noL1= is prescaled, the \verb=HLT_2mu4= decision is used.



Once the single-muon efficiencies are known, the efficiencies given
  in Eq.~\eqref{eq:2mu3_eff} and Eq.~\eqref{eq:mu4_mu4_nol1_eff} can be calculated.
The evaluation of the single-muon trigger efficiencies is described below.


\subsubsection{Triggers efficiencies for HLT\_mu4 Trigger}
The single muon efficiencies for the \verb=HLT_mu4= trigger are measured using the \PbPb\ minimum bias
  data samples: events recorded in the min-bias stream satisfying either one
  of the following triggers (Described in Section~\ref{sec:data}):
\begin{itemize}
\item \verb=HLT_noalg_cc_L1TE600_0ETA49=
\item \verb=HLT_noalg_pc_L1TE50_VTE600_0ETA49=
\item \verb=HLT_mb_sptrk_L1ZDC_A_C_VTE50=.
\end{itemize}
The single-muon trigger efficiency for \verb=HLT_mu4= 
  is defined to be the fraction of reconstructed (Tight or Medium) muons
  that have a matching HLT muon within $\Delta R < 0.1$:
\begin{eqnarray}
\varepsilon_{\mathrm{trig}}^{\mathrm{mu4}}=
\frac{\text{Number of reconstructed muons in MinBais Events matching the HLT\_mu4 trigger}}{\text{Number of reconstructed muons in MinBias Events}},
\label{eq:trig_eff_mb}
\end{eqnarray}
  where the ratio in Eq.~\ref{eq:trig_eff_mb} are evaluated as a function of \pt\ and $q*\eta$.
The efficiencies, for muons passing the Tight working point,
  measured as a function of \pT\ and $q*\eta$,
  are shown in Figure~\ref{fig:pbpb_trigEff} as a 2D plot, and in
  Figure~\ref{fig:pbpb_trigEff_1D2}
  in different intervals of $q*\eta$.
Appendix~\ref{sec:medium_muon_trigEff} shows similar 
  efficiencies but for muons passing the Medium working point.
For \pt>6~\GeV, the efficiencies at forward (or backward) $q*\eta$ are larger
  than at mid-rapidity.
While at low-pt the efficiencies decrease from negative $q*\eta$ to positive
  $q*\eta$; this is most clearly seen in the 2D plots shown in Figure~\ref{fig:pbpb_trigEff}.

The lines in Figure~\ref{fig:pbpb_trigEff_1D2}
  indicate fits to Fermi functions + linear function of the form:
\begin{eqnarray}
\mathrm{eff}(\pt)=\frac{\epsilon}{1+\exp{(\frac{\pt-p_{\mathrm{T}}^0}{\Delta})}} + A\times\Theta(\pt-p_{\mathrm{T}}^{\mathrm{Thres}})(\pt-p_{\mathrm{T}}^{\mathrm{Thres}})
\label{eq:trigger_eff_fit}
\end{eqnarray}
where $\epsilon$, $p_{\mathrm{T}}^0$, $\Delta$ and $A$ are fit-parameters,
  $\Theta$ is the step function,
  and $p_{\mathrm{T}}^{\mathrm{Thres}}$ is 4~\GeV\ for the \verb=HLT_mu4= trigger.
The fits are made separately for each $q*\eta$ interval and the fit parameters
  vary as function of $q*\eta$.
The linear function is included in the parameterization to account
  for an increasing trend observed in the efficiency in few
  $q*\eta$ intervals.
In case the fit function in Eq.~\eqref{eq:trigger_eff_fit} 
  returns a value greater than 1, which can happen at very high-\pt\ 
  due to the presence of the linear term, it is trunated to 1.0. 
Also, the fits are not perfect at low-\pt in a few $q*\eta$ intervals.
However, since the data-points themselves have very small errors,
  for $\pt<8$~\GeV\ the data-points are interpolated to obtain the efficiencies
  and the fits are used only for $\pt>8$~\GeV.
The trigger efficiencies show a few percent variation with centrality,
  this is discussed later in Sec~\ref{sec:trigeff_centdep}.
%This is demonstrated in Figure~\ref{fig:pbpb_trigEff_1D2_centdep}.
%The efficiency corrections are thus applied as function of 
%  centrality,\pt, and $q*\eta$.




\begin{figure} 
\centering
%\includegraphics[width=0.5\textwidth]%
%{figures/AllFigs/Can_trigeff_type1_trig0_MuonSelectionCut0}%
\includegraphics[width=0.5\textwidth]%
{figures/AllFigs/Can_trigeff_type1_trig0_MuonSelectionCut1}
\caption{
HLT\_mu4 trigger efficiency for \PbPb.
The Plot corresponds to Tight muons.
\label{fig:pbpb_trigEff}
}
\end{figure}



\begin{figure} \centering
\includegraphics[width=1\textwidth]%
{figures/AllFigs/Can_eff_trig_option0_MuonSelectionCut1}%
\caption{
HLT\_mu4 trigger efficiency for \PbPb.
The lines indicate fits to Fermi+linear functions given in Eq.~\eqref{eq:trigger_eff_fit}.
The fits are made over the range \pt>4~\GeV.
The plots correspond to Tight muons.
\label{fig:pbpb_trigEff_1D2}
}
\end{figure}



\begin{comment}
\begin{figure} \centering
\includegraphics[width=0.5\textwidth]%
{figures/AllFigs/Can_eff_cent2_trig0_ptmin40_ptmax200_option0_MuonSelectionCut1}%
\caption{
The relate change in the trigger efficiency for 
  tight muons as a function of centrality (averaged over $\eta$).
\label{fig:pbpb_trigEff_1D2_centdep}
}
\end{figure}
\end{comment}



\FloatBarrier
\clearpage
\subsubsection{Triggers efficiencies for HLT\_mu4noL1 and  HLT\_mu4noL1$\cap$HLT\_mu4 Triggers}
The efficiency for the \verb=HLT_mu4_mu4noL1= trigger and the combined
  efficiency for \verb=HLT_mu4noL1=$\cap$\verb=HLT_mu4= are obtained using 
  the HardProbes stream, by using events that are triggered by 
  the following single muon triggers.
\begin{itemize}
\item \verb=HLT_mu4=
\item \verb=HLT_mu6=
\item \verb=HLT_mu8=.
\end{itemize}
Offline muon-pairs are selected where one of the muon matches
  the \verb=HLT_mu4= or \verb=HLT_mu6= or \verb=HLT_mu8= triggers 
  (we will call this muon as \textit{first muon} in the text below).
Then the fraction of such offline muon pairs that match the 
  \verb=HLT_mu4_mu4noL1= trigger (or simultaneously match both \verb=HLT_mu4_mu4noL1= and \verb=HLT_2mu4=) 
  are evaluated as a function of the \pT\ and $q\times\eta$ of the other muon
  in the pair (this muon is called \textit{second muon} in the text below).
This fraction is the efficiency of the \verb=HLT_mu4noL1= trigger (or \verb=HLT_mu4noL1=$\cap$\verb=HLT_mu4= Trigger).
For example, the \verb=HLT_mu4noL1= can be obtained as:
\begin{eqnarray}
\varepsilon_{\mathrm{trig}}^{\mathrm{mu4\_noL1}}&=&
\frac{N^{\mathrm{pairs}}\{\text{Matching HLT\_mu4\_mu4noL1, where first muon in pair matches HLT\_mu4}\}}{N^{\mathrm{pairs}}\{\text{First muon in pair matches HLT\_mu4 trigger}\}}, \\
  &=& \text{Prob}(\text{HLT\_mu4\_mu4noL1}(\vec{p}_{a},\vec{p}_{b})|\text{HLT\_mu4}(\vec{p}_{a})) \nonumber
\label{eq:trig_eff_hp}
\end{eqnarray}
where, the efficiency of the l.h.s. of the above equation is 
  evaluated as function of \pt\ and $q*\eta$ of the second muon in the pair.
The efficiencies for the \verb=HLT_mu4noL1= trigger 
  measured as a function of \pT\ and $q\times\eta$,
  are shown in Figure~\ref{fig:pbpb_trigEff_noL1}, as 2D plots and in
  Figures~\ref{fig:pbpb_trigEff_1D1_noL1} in different intervals of $q\times\eta$.
Figure~\ref{fig:pbpb_trigEff_1D1_noL1} also shows the 
  combined \verb=HLT_mu4noL1=$\cap$\verb=HLT_mu4= efficiency.
Figure~\ref{fig:pbpb_trigEff_1D1_noL1} also shows the efficiency 
  for the HLT\_mu4 trigger obtained using the HP stream.

As a check, the \verb=HLT_mu4= efficiency obtained from the HP stream is compared
  with the efficiency obtained from the min-bias stream
  in Figure~\ref{fig:pbpb_trigEff_1D1_compare}.
The two estimates of the efficiency are found to be consistent.
This plot demonstrates that the method used here to obtain the 
  single-muon efficiency from the HP stream works (otherwise the
  estimate from Eq.~\eqref{eq:trig_eff_hp} would not match the estimate
  from the MinBias events).
This plot also implicitly verifies that the facorization
  of the trigger efficiency, Eq.\eqref{eq:2mu3_eff}, is correct.
As it shows that :
\begin{eqnarray}
\text{Prob}(\text{HLT\_mu4}(\vec{p}_{b})) = \text{Prob}(\text{HLT\_2mu4}(\vec{p}_{a},\vec{p}_{b})|\text{HLT\_mu4}(\vec{p}_{a}))
\label{eq:2mu4_factorization}
\end{eqnarray}
Appendix~\ref{sec:medium_muon_trigEff} shows similar 
  efficiencies but for muons passing the Medium working point.

\begin{figure}
\centering
\includegraphics[width=0.5\textwidth]%
{figures/AllFigs/can_eff_Mu4NoL1_RefAll_MuonSelectionCut1}%
\includegraphics[width=0.5\textwidth]%
{figures/AllFigs/can_eff_Both_RefAll_MuonSelectionCut1}
\caption{
  Left: HLT\_mu4noL1 trigger efficiency for \PbPb.
  Right: HLT\_mu4noL1$\cap$HLT\_mu4 Triggers trigger efficiency for \PbPb.
  Plots correspond Tight muons.
\label{fig:pbpb_trigEff_noL1}
}
\end{figure}


\begin{figure}
\centering
\includegraphics[width=1\textwidth]%
{figures/AllFigs/can_eff1D_All_RefAll_MuonSelectionCut1}%
\caption{
  Efficiencies for the HLT\_mu4NoL1 trigger (red) as well as the combined 
    HLT\_mu4noL1$\cap$HLT\_mu4 eficiency (blue, labelled as \textit{Both}) in \PbPb.
  Each panel corresponds to a different $q\eta$ interval.
  The efficiency for the HLT\_mu4 trigger obtained using the HP
    stream is also shown (black).
  The plots correspond to Tight muons.
\label{fig:pbpb_trigEff_1D1_noL1}
}
\end{figure}



\begin{figure}
\centering
\includegraphics[width=1\textwidth]%
{figures/AllFigs/can_eff1D_Mu4_RefAll_MuonSelectionCut1}%
\caption{
  Comparison of the efficiency for the HLT\_mu4 trigger obtained using the HP
    stream to that using the min-bias events.
  The plots correspond to Tight muons.
  This plot also demonstrates the validity of the factorization~Eq.~\eqref{eq:2mu3_eff},
    via Eq.~\eqref{eq:2mu4_factorization}.
\label{fig:pbpb_trigEff_1D1_compare}
}
\end{figure}



\FloatBarrier
\subsubsection{Triggers efficiencies: Centrality dependence~\label{sec:trigeff_centdep}}
The trigger efficiencies also have some centrality dependence.
This is demonstrated in Figure~\ref{fig:pbpb_trigEff_centdep}.
The plot shows the relative variation in the efficiencies w.r.t. 
  the centrality-integrated efficiency, averaged over all $q\eta$ intervals.
This additional centrality dependence is included when applying the 
  trigger efficiency corrections.
This factor is only implemented as a function of centrality and \pt, as there are 
  not enough statistics to implement it as function of (centrality,\pt,$q\eta$).
The centrality dependence is evaluated using tight muons as they are less susceptible
  to misidentified muons.
The final efficiency (for a given trigger) is obtained by taking the product 
  of this centrality-dependent factor (Figure~\ref{fig:pbpb_trigEff_centdep})
  and the centrality averaged efficiency (for ex. Figure~\ref{fig:pbpb_trigEff_1D1_noL1}):
\begin{align}
	\epsilon_{\mathrm{trig}}(\pt,q\eta,\mathrm{centrality}) = \epsilon_{\mathrm{trig}}(\pt,q\eta)\times\mathrm{CentDep}(\pt,\mathrm{centrality})
\label{eq:trigeff_centdep}
\end{align}

\begin{figure}
\centering
\includegraphics[bb=382bp  0bp 565bp 138bp,clip,width=0.33\textwidth]%
{figures/AllFigs/can_eff_centdep_Mu4_RefAll40_50_MuonSelectionCut1}%
\includegraphics[bb=382bp  0bp 565bp 138bp,clip,width=0.33\textwidth]%
{figures/AllFigs/can_eff_centdep_Mu4NoL1_RefAll40_50_MuonSelectionCut1}%
\includegraphics[bb=382bp  0bp 565bp 138bp,clip,width=0.33\textwidth]%
{figures/AllFigs/can_eff_centdep_Both_RefAll40_50_MuonSelectionCut1}
\includegraphics[bb=382bp  0bp 565bp 138bp,clip,width=0.33\textwidth]%
{figures/AllFigs/can_eff_centdep_Mu4_RefAll50_60_MuonSelectionCut1}%
\includegraphics[bb=382bp  0bp 565bp 138bp,clip,width=0.33\textwidth]%
{figures/AllFigs/can_eff_centdep_Mu4NoL1_RefAll50_60_MuonSelectionCut1}%
\includegraphics[bb=382bp  0bp 565bp 138bp,clip,width=0.33\textwidth]%
{figures/AllFigs/can_eff_centdep_Both_RefAll50_60_MuonSelectionCut1}
\includegraphics[bb=382bp  0bp 565bp 138bp,clip,width=0.33\textwidth]%
{figures/AllFigs/can_eff_centdep_Mu4_RefAll60_80_MuonSelectionCut1}%
\includegraphics[bb=382bp  0bp 565bp 138bp,clip,width=0.33\textwidth]%
{figures/AllFigs/can_eff_centdep_Mu4NoL1_RefAll60_80_MuonSelectionCut1}%
\includegraphics[bb=382bp  0bp 565bp 138bp,clip,width=0.33\textwidth]%
{figures/AllFigs/can_eff_centdep_Both_RefAll60_80_MuonSelectionCut1}
\includegraphics[bb=382bp  0bp 565bp 138bp,clip,width=0.33\textwidth]%
{figures/AllFigs/can_eff_centdep_Mu4_RefAll80_200_MuonSelectionCut1}%
\includegraphics[bb=382bp  0bp 565bp 138bp,clip,width=0.33\textwidth]%
{figures/AllFigs/can_eff_centdep_Mu4NoL1_RefAll80_200_MuonSelectionCut1}%
\includegraphics[bb=382bp  0bp 565bp 138bp,clip,width=0.33\textwidth]%
{figures/AllFigs/can_eff_centdep_Both_RefAll80_200_MuonSelectionCut1}
\caption{
The centrality dependence of the single-muon trigger efficiencies.
The plot shows the relative variation in the efficiencies w.r.t. 
  the cenrality-integrated efficiency, averaged over all $q\eta$ intervals.
The left, center and right panels correspond to the HLT\_mu4, HLT\_mu4noL1 and
  HLT\_mu4noL1$\cap$HLT\_mu4 (labelled as ``both'') triggers respectively.
From top to bottom the rows correspond to different muon-\pt\ intervals.
The red-line shows the parameterization of the centrality dependence 
  that is used in the analysis, the blue lines show a $\pm$3\% band
  for reference %which is included as an uncertainty.
\label{fig:pbpb_trigEff_centdep}
}
\end{figure}






\clearpage
\subsection{Reconstruction efficiencies\label{sec:reco_eff}}
The single muon reconstruction efficiency is evaluated using the \starlight\
  MC samples~\cite{ATLHI-304,ATLHI-313} with HIJING Overlay 
  discussed in Section~\ref{sec:mc_samples}.
The HIJING samples with 2015 and 2018 conditions show only a 
  minor performance difference and the two samples are combined 
  while evaluating the efficiencies.
The efficiency is defined as the fraction of the number of generated-muons that have
  associated reconstructed muons -- obtained with the same requirements as those
  used in the data analysis -- that match the generated-muon with a probability
  of 0.5 or more.
\begin{align}
\epsilon_{\mathrm{reco}}(\pt,q\eta,\mathrm{centrality}) =\frac{N^{\mathrm{Gen}}(\pt,q\eta,\mathrm{centrality})\text{ with match-prob>0.5 to reco muon passing selections}}{N^{\mathrm{Gen}}(\pt,q\eta,\mathrm{centrality})}
\end{align}
Figure~\ref{fig:pbpb_ReconEff_2D2} shows the reconstruction efficiency
  as a function of \pt and $q*\eta$.
Figure~\ref{fig:pbpb_ReconEff_tight} shows the reconstruction efficiency
  a function of \pT\ for several $q*\eta$ bins.
%\begin{eqnarray}
%  \mathrm{eff}(\pt)=\frac{\epsilon}{1+\exp{(\frac{\pt-\pt^0}{\Delta})}}
%\label{eq:reco_eff_fit}
%\end{eqnarray}
%where $\epsilon$, $\pt^0$ and $\Delta$ are fit-parameters.
Appendix~\ref{sec:medium_muon_recoeff} shows similar 
  efficiencies but for muons passing the Medium working point.
%In general the Fermi-function fits do not work well in many regions 
%  for the Tight muons.
The evaluated points themselves have quite small uncertaintues,
  so instead of fiting efficiencies with fermi functions, 
  the evaluated efficiency data-points are interpolated and 
  used in the analysis.

For brevity, the efficiencies shown in Figure~\ref{fig:pbpb_ReconEff_tight} 
  are without any centrality binning i.e. they are from all events in 
  the Overlay MC sample.
However in the analysis the efficiencies are evaluated in bins of 
  (10\% or 20\% wide) centrality.
In general the efficiencies show a small but systematic dependence 
  on centrality and increase from central to peripheral events.
Thus the efficiency corrections are applied as a function of 
  \pt, $q*\eta$ and centrality.
These efficiencies, in narrow bins of centrality, 
  are shown in Appendix~\ref{sec:recoeff_centdep}.
On top of these MC-driven reconstruction efficiency estimates,
  additional scale-factors as recommended by the MuonCP
  group~\cite{MCPAnalysisWinterMC16}, and provided by the
  CP::MuonEfficiencyScaleFactors tool~\cite{MuonEfficiencyScaleFactors}
  are included.

\begin{figure} 
\centering
%\includegraphics[width=0.5\textwidth]%
%{figures/AllFigs/Can_recoeff2D_medium}%
\includegraphics[width=0.5\textwidth]%
{figures/AllFigs/Can_recoeff2D_tight}
\caption{
Muon reconstruction efficiencies for \PbPb.
The plot corresponds to Tight muons.
\label{fig:pbpb_ReconEff_2D2}
}
\end{figure}




\begin{figure}
\centering
\includegraphics[width=1.0\textwidth]%
{figures/AllFigs//Can_recoeff1D_tight}
\caption{
Muon reconstruction efficiency as a function of \pt and
  for several $q*\eta$ slices.
Plots are for Tight quality muons.
\label{fig:pbpb_ReconEff_tight}
}
\end{figure}


\FloatBarrier
Figure~\ref{fig:check_pair_effs} shows the net efficiency correction 
  (Eq.~\eqref{eq:net_eff_correction}) for pairs that pass 
  the selections described later in Section~\ref{sec:analysis}.
\begin{figure}
\centering
\includegraphics[width=1.0\textwidth]%
{figures/AllFigs/can_Eff_Corrections_pTBar_GT5_EffCor}
\includegraphics[width=1.0\textwidth]%
{figures/AllFigs/can_Eff_Corrections_Tight_pTBar_GT5_EffCor}
\caption{
Left : the spectra for number of pairs as a function of the 
  \pt\ of the lower-\pt\ muon in the pair.
Middle: the mean combined efficiency correction,$\langle w\rangle$ (Eq.~\eqref{eq:net_eff_correction})
  on the pair as a function
  of the \pt\ of the lower-\pt\ muon in the pair.
Right: the distribution of the efficiency correction.
Plots are for \ptab>4~\GeV and \ptbar>5~\GeV,
  where \ptbar\ is the average \pt\ of the two muons in the pair.
The top panels are for Medium muons and the lower ones are for Tight muons.
\label{fig:check_pair_effs}
}
\end{figure}

\clearpage


